% Canonical Paper II source.
Paper I develops Arithmetic Expression Geometry (AEG) from ordered evaluation
histories to affine flow, regular expression surfaces, and contact curvature
\cite{Yuan2026AEGFoundations}.  The present paper asks the next question: once the
geometric data have been fixed, what analysis do they actually support?  We answer it
first on a general regular arithmetic expression space and then completely enough on
the basic hyperbolic model to obtain a boundary kernel, a Dirichlet theorem, and an
exact boundary energy operator.

This question has two layers which must not be conflated.  A regular AES
\[
  \mathfrak E=(M,g,a;\mu,\lambda)
\]
is an oriented Riemannian \emph{surface}.  Its metric and orientation already
determine rotation by $+\pi/2$, hence the usual complex structure of an oriented
Riemannian surface.  The assignment $a$ selects a distinguished positively oriented
orthonormal frame $(X_u,X_v)$.  Complex and elliptic analysis on this surface are
therefore intrinsic to the data imported from Paper I.

The contact state space is instead
\[
 \mathcal C=\R^3_{(u,v,a)},\qquad
 \alpha=da-\mu\,du-\lambda a\,dv,
\]
with preferred horizontal lifts
\[
 D_u=\partial_u+\mu\partial_a,
 \qquad
 D_v=\partial_v+\lambda a\partial_a.
\]
When $\mu\lambda\ne0$, its horizontal distribution is nonintegrable.  To do
analysis there one must additionally choose a horizontal metric, its compatible
complex structure, a positive measure, and operator domains.  The resulting object
is a three-dimensional contact-CR problem, not another presentation of the
two-dimensional surface problem.

Figure~\ref{fig:two-analytic-layers} records this separation.  Both branches retain
the arithmetic derivative rules
\[
 X_ua=\mu,\quad X_va=\lambda a,
 \qquad
 D_ua=\mu,\quad D_va=\lambda a,
\]
but their brackets have different geometric meanings.  On a surface the bracket is
tangent.  On $\mathcal C$ one has
$[D_u,D_v]=\mu\lambda\partial_a$, a vertical curvature term.

\begin{figure}[t]
\centering
\resizebox{0.95\textwidth}{!}{%
\begin{tikzpicture}[
  node distance=8mm and 8mm,
  box/.style={draw,rounded corners,align=center,minimum width=4.3cm,
    minimum height=1.05cm,fill=blue!3},
  result/.style={draw,rounded corners,align=center,minimum width=4.3cm,
    minimum height=1.05cm,fill=green!5},
  arr/.style={-{Latex[length=2mm]},thick},
  warn/.style={font=\scriptsize,align=center,text=red!60!black}
]
  \node[box] (paperi) {Paper I affine propagation\\and assignment derivative rules};
  \node[box,below left=of paperi] (surface) {regular AES surface\\$(M,g,a;\mu,\lambda)$};
  \node[box,below right=of paperi] (contact) {contact state space\\$(\mathcal C,\alpha,D_u,D_v)$};
  \node[result,below=of surface] (elliptic) {canonical surface\\complex structure\\Laplace--Beltrami and Poisson theory};
  \node[result,below=of contact] (cr) {chosen horizontal\\metric and measure\\CR and sub-Laplacian theory};
  \draw[arr] (paperi)--(surface);
  \draw[arr] (paperi)--(contact);
  \draw[arr] (surface)--(elliptic);
  \draw[arr] (contact)--(cr);
  \draw[<->,very thick,dashed,gray!65] (elliptic)--(cr)
    node[midway,above,warn] {structural comparison\\not identification};
\end{tikzpicture}%
}
\caption{The two analytic objects carried forward from Paper I.  The surface branch
is elliptic and has an intrinsic complex structure.  The contact branch is
subelliptic and requires an analytic normalization.  Nonintegrability prevents the
contact distribution from being identified with a regular AES surface.}
\label{fig:two-analytic-layers}
\end{figure}

\subsection{Principal results}

Our first result fixes the real operator theory on a regular AES.  If
\[
 [X_u,X_v]=c_uX_u+c_vX_v,
\]
then, with the divergence-of-gradient convention,
\begin{equation}\label{eq:intro-surface-laplacian}
 \Delta_g=X_u^2+X_v^2-c_vX_u+c_uX_v.
\end{equation}
The first-order drift is forced by the Riemannian volume and cannot be discarded.
The compatible Cauchy--Riemann operator
\[
 \dbarAEG=\frac12(X_u+iX_v)
\]
then satisfies an adjoint factorization of $-\Delta_g$.  In particular,
arithmetic holomorphic fields are genuinely Laplace--Beltrami harmonic.  We also
show that $\dbarAEG$ differs from the ordinary Riemann-surface operator only by a
nowhere-zero unitary gauge factor.  Thus the surface theory does not manufacture a
new class of holomorphic functions: its AEG content is the assignment-adapted frame,
the distinguished solution families, and their arithmetic interpretation.

Our second result fixes the corresponding ambiguity on the contact state space.
We normalize the horizontal metric by declaring $(D_u,D_v)$ orthonormal and use the
positive contact volume.  On compactly supported smooth fields,
\begin{equation}\label{eq:intro-contact-adjoints}
 D_u^*=-D_u,
 \qquad
 D_v^*=-D_v-\lambda.
\end{equation}
Consequently the variational sub-Laplacian is
\begin{equation}\label{eq:intro-contact-laplacian}
 \Delta_{\mathcal C}
 =D_u^2+D_v^2+\lambda D_v,
 \qquad
 -\Delta_{\mathcal C}=D_u^*D_u+D_v^*D_v.
\end{equation}
The legacy expression $D_u^2+D_v^2$ remains useful, but only as the raw sum of
squares occurring in the formal CR factorization.  We prove both the raw and adjoint
factorizations, relate two natural measures by a unitary gauge transformation, and
give a Friedrichs realization of the nonnegative energy operator.  H\"ormander's
bracket criterion supplies hypoellipticity of the differential expression
\cite{Hormander1967Hypoelliptic}.

Between the local operator theory and the global hyperbolic boundary problem, we
establish a functorial construction that will be used by the singular theory.  The
planar data
\[
 A_\phi(w)=\operatorname{Im}(e^{-i\phi}w),
 \qquad
 h_\phi=\frac{|dw|^2}{\mu^2+\lambda^2A_\phi(w)^2}
\]
form a regular AES with harmonic assignment.  Every locally biholomorphic map
\(F\) pulls this data back to a regular AES; its Gaussian curvature is
\[
 K=\lambda^2\frac{\mu^2-\lambda^2a^2}{\mu^2+\lambda^2a^2}.
\]
We also construct the complete flat model on \(\C^*\),
\[
 s=\log|W|,\qquad
 a=\frac{\mu}{\lambda}\sinh(\lambda s),
 \qquad
 g=\left|\frac{dW}{W}\right|^2,
\]
with the continuous limiting formula \(a=\mu s\) when \(\lambda=0\).  Its zero
set is the unit circle, and a locally biholomorphic pullback has zero set
\(W^{-1}(S^1)\).  A descent proposition covers unitary automorphy and the possible
sign character arising from inversion.  The construction is deliberately stated
only away from critical points and branch points; the cylindrical construction is
in addition restricted away from zeros and poles of $W$.

Our decisive global theorem concerns the basic hyperbolic AES.  For
$\mu,\lambda>0$, put
\[
 X=\frac{\lambda}{\mu}x,
 \qquad
 w=X+iy.
\]
Then
\[
 g_{\mu,\lambda}
 =\frac1{\lambda^2}\frac{dX^2+dy^2}{y^2},
 \qquad
 \Delta_g=\lambda^2y^2(\partial_X^2+\partial_y^2),
\]
and the arithmetic holomorphic equation is precisely
$\partial_{\bar w}F=0$.  The standard upper-half-plane Poisson kernel therefore
transports exactly to the AEG normalization.  For every
$\varphi\in C_0(\R)$ we prove existence and uniqueness in the class
$C^2(\HH)\cap C(\overline\HH^{\,\mathrm{conf}})$ with the prescribed value zero at
$\infty$.  For Schwartz data,
the conformal-boundary Dirichlet-to-Neumann map is
\[
 \mathcal N=|D_X|,
\]
and the hyperbolic Dirichlet energy is exactly its homogeneous
$H^{1/2}$ quadratic form.  This is the first exact boundary-operator statement of
the AEG analytic programme.  Its classical analytic ingredients are
standard \cite{Ahlfors1979ComplexAnalysis,CaffarelliSilvestre2007Extension}; the
result here is their explicit transport, domain statement, and assignment-coordinate
interpretation on the Paper I model.

The same model supplies several genuinely assignment-dependent fields because
\begin{equation}\label{eq:intro-w-assignment}
 w=y\left(i-\frac{\lambda}{\mu}a\right).
\end{equation}
Thus every holomorphic $\Phi(w)$ is an explicit field in $(a,y)$ coordinates.
We also classify all scalar harmonic fields depending only on $a$:
\[
 h(a)=C_0+C_1\arctan\!\left(\frac{\lambda a}{\mu}\right).
\]
The nonconstant member is the harmonic measure of a boundary half-line.  In the
contact branch, logarithmic CR first integrals and finite filtered affine--Appell
modules provide a separate collection of assignment-dependent solutions.

\subsection{Boundaries of the claim}

Several restrictions are essential.

\begin{itemize}
\item The line $y=0$ is the conformal or ideal boundary of the complete hyperbolic
metric, not a finite Riemannian boundary component.
\item Poisson extensions with nonzero ideal-boundary values are generally not in
$L^2(\HH,d\vol_g)$.  The boundary-energy theorem and the $L^2$ Friedrichs
realization therefore occupy different analytic domains.
\item The identity $\Delta_ga=2\lambda^2a$ from Paper I is pointwise.  The assignment
$a=-x/y$ is not an $L^2$ eigenvector.
\item We do not claim that the normalized contact CR structure is invariant under all
contactomorphisms, nor that its explicit CR families are complete in a Hilbert or
Banach space.
\item We do not establish Green kernels, spectral completeness, or continuation on a
general regular or singular AES.
\item Holomorphic pullbacks are regular AES constructions only where the map is
locally biholomorphic.  For the cylindrical target, its coordinate $W$ must in
addition be finite and nonzero.  Critical points, poles, algebraic branch points,
and singular zero-graph completion belong to Paper III.
\end{itemize}

These limitations are not peripheral.  They separate proved function theory from
the open analytic programme and keep the interfaces to singular geometry and
projective complexity acyclic.

\subsection{Organization}

Section~\ref{sec:analytic-data} fixes measures, domains, structure coefficients, and
energies on a regular AES.  Section~\ref{sec:surface-operators} develops the
intrinsic surface Cauchy--Riemann theory and proves holomorphicity implies
Laplace--Beltrami harmonicity.  Section~\ref{sec:contact-analysis} treats the
distinct contact-CR problem.  Section~\ref{sec:pullback-cylinder} proves the planar
pullback theorem, constructs the complete cylindrical AES, and states the regular
automorphic descent interface.  Section~\ref{sec:hyperbolic-analysis} identifies
all operators on the basic model.  Sections~\ref{sec:poisson-dirichlet} and
\ref{sec:boundary-energy} prove the boundary theorems and construct explicit
assignment-dependent families.  The appendices contain the full frame, contact, and
Fourier calculations.
